\subsection{微分形式，外微分}

回顾一下，在本篇以及后续各篇中，我们假定 K 的特征为 0，或者所考虑的簇在局部具有有限维。

8.3.1. 设 X 是 \(\mathbf{C}^r\) 类流形，且 F 是以 X 为基底的 \(\mathbf{C}^k\) 类向量丛，其中 \(k + 1 \in \mathbb{N}_{\mathbb{K}}\) 且 \(k + 1 \leqslant r\)。设 \(p\) 是一个 \(\geqslant 0\) 的整数，并令 \(\operatorname{Alt}^p(\mathrm{T}(\mathrm{X}); \mathrm{F})\) 为一个 \(\mathbf{C}^k\) 类向量丛，它由 \(p\)-交错多线性映射（从 \(\mathrm{T}(\mathrm{X})\) 到 F）组成（7.8.1 及第 7.8 号勘误）。X 的开集 U 上的这个丛的一个截面称为 U 上取值于 F、次数为 \(p\) 的交错微分形式（或简称微分形式）。属于 \(\mathbf{C}^k\) 类的这些截面构成一个 \(\mathcal{C}^k(\mathrm{U})\)-模，记为 \(^k\Omega^p(\mathrm{U}; \mathrm{F})\)。若 E 是 Banach 空间，则记 \(^k\Omega^p(\mathrm{U}; \mathrm{E})\) 而不记 \(^k\Omega^p(\mathrm{U}; \mathrm{E}_{\mathrm{x}})\)，并称为取值于 E 的微分形式。在这些记号中，有时省略 \(k\)（分别为 E），当它等于 \(r - 1\)（分别为 K）时。

取标量值的 0 次（分别为 1 次）微分形式是函数（分别是协向量场）。

设 \(\varphi\colon\mathbf{Y}\to\mathbf{X}\) 是 \(\mathbf{C}^{r}\) 类流形的态射，且 \(\omega\in{}^{k}\Omega^{p}(\mathbf{X};\mathbf{F})\)。则存在且仅存在一个形式 \(\varphi^{*}\omega\in{}^{k}\Omega^{p}(\mathbf{Y};\varphi^{*}\mathbf{F})\)，使得

\[
(\varphi^ {*} \omega) _ {y} (v _ {1}, \dots , v _ {p}) = \omega_ {\varphi (y)} (\mathrm{T} _ {y} (\varphi). v _ {1}, \dots , \mathrm{T} _ {y} (\varphi). v _ {p})\tag{1}
\]

对所有 \(y \in Y\) 以及 \(v_{1}, \ldots, v_{p}\) 中元素的任意族 \(\mathrm{T}_{y}(\mathrm{Y})\) 都成立。

若 \(\psi: Z \to Y\) 是 \(C^r\) 类流形的态射，则有

\[
(\varphi \circ \psi) ^ {*} \omega = \psi^ {*} (\varphi^ {*} \omega).\tag{2}
\]

当 Y 是 X 的子流形且 \(\varphi\) 是典范嵌入时，有时用 \(\omega|Y\) 代替 \(\varphi^{*}\omega\)，并称 \(\omega|Y\) 是由 \(\omega\) 在 Y 上诱导的形式。

8.3.2. 第 7.8 号的定义和结果适用于微分形式。特别地，向量丛的一个配对 \(F' \times_{\mathbf{X}} F'' \to F\) 产生一个外积 \(^k\Omega^{p'}(\mathrm{U};\mathrm{F}') \times ^k\Omega^{p''}(\mathrm{U};\mathrm{F}'') \to ^k\Omega^p(\mathrm{U};\mathrm{F})\)，其中 \(p = p' + p''\)（7.8.2）；记作 \((\omega', \omega'') \mapsto \omega' \wedge \omega''\)。因此，对所有 \(^k\Omega^p(\mathrm{U};\mathrm{K})\) 的 \(p \geqslant 0\) 的直和带有结合且交错的分次代数结构（A, III, p. 53）。

设 \(\xi\) 是 X 上的向量场，\(\omega\) 是 X 上取值于向量丛 F、次数为 \(p \geqslant 1\) 的微分形式；\(\xi\) 与 \(\omega\) 的内积是 X 上取值于 F、次数为 \(i(\xi, \omega)\) 的微分形式 \(p - 1\)，定义为

\[
i (\xi , \omega) _ {x} (v _ {1}, \dots , v _ {p - 1}) = \omega_ {x} (\xi (x), v _ {1}, \dots , v _ {p - 1})\tag{3}
\]

对所有 \(x \in X\) 以及 \(v_{1}, \ldots, v_{p-1}\) 中元素的任意族 \(\mathrm{T}_{x}(\mathrm{X})\) 都成立。若 \(\omega\) 是 0 次微分形式，令 \(i(\xi, \omega) = 0\)。我们也写作 \(i(\xi)\omega\) 或 \(i_{\xi}\omega\)，以代替 \(i(\xi, \omega)\)。当 \(\mathrm{F} = \mathrm{K}_{\mathrm{X}}\) 时，上述定义与 7.8.4 的定义一致。

设 \(\mathbf{F}' \times_{\mathbf{x}} \mathbf{F}'' \to \mathbf{F}\) 是向量丛的一个配对。对于 \(\omega' \in {}^k \Omega^{p'}(\mathbf{U}; \mathbf{F}')\) 和 \(\omega'' \in {}^k \Omega^{p''}(\mathbf{U}; \mathbf{F}'')\)，有

\[
i (\xi) \left(\omega^ {\prime} \wedge \omega^ {\prime \prime}\right) = \left(i (\xi) \omega^ {\prime}\right) \wedge \omega^ {\prime \prime} + (- 1) ^ {p ^ {\prime}} \omega^ {\prime} \wedge i (\xi) \omega^ {\prime \prime}.\tag{4}
\]

8.3.3. 设 \((u^{1},\ldots,u^{n})\) 是 X 的开集 U 上的坐标系。\(\omega\) 中的每个元素 \({}^{k}\Omega^{p}(\mathrm{U};\mathrm{F})\) 都唯一地写成

\[
\omega = \sum_ {i _ {1} <   \dots <   i _ {p}} f _ {i _ {1}, \dots , i _ {p}}. d u ^ {i _ {1}} \wedge \dots \wedge d u ^ {i _ {p}}
\]

其中 \(f_{i_1,\ldots,i_p}\) 是 \(\mathbf{C}^k\) 上 \(\mathbf{F}\) 的 \(\mathbf{U}\) 类截面；在此公式中，符号 \(\wedge\) 所涉及的是典范配对 \(\mathbf{K}_{\mathbf{x}} \times_{\mathbf{x}} \mathbf{K}_{\mathbf{x}} \to \mathbf{K}_{\mathbf{x}}\)，而表示乘积的点号所涉及的是典范配对 \(\mathbf{F} \times_{\mathbf{x}} \mathbf{K}_{\mathbf{x}} \to \mathbf{F}\)。

8.3.4. 设 \(p\) 是一个 \(\geqslant 0\) 的整数，\(r\) 是 \(\mathbf{N}_{\kappa}\) 的一个元素。设 \(E\) 和 \(F\) 是 Banach 空间，\(U\) 是 \(E\) 的开子集，且 \(\alpha \in {}^r\Omega^p(U;F)\)。令 \(\tilde{\alpha} \in \mathcal{C}^r (U; \operatorname{Alt}^p(E;F))\) 为对应于 \(\alpha\) 的元素。若 \(x \in U\)，则 \(d_x\tilde{\alpha}\) 是 \(\tilde{\alpha}\) 在 \(x\) 处的微分（5.5.6），它属于 \(\mathcal{L}(E; \operatorname{Alt}^p(E;F))\)；若 \(t \in E\)，则 \((d_x\tilde{\alpha})(t) \in \operatorname{Alt}^p(E;F)\)；而若 \(t_1, \ldots, t_p\) 属于 \(E\)，则 \((d_x\tilde{\alpha})(t)(t_1, \ldots, t_p) \in F\)。存在且仅存在一个元素 \(d\alpha\)；它属于 \({}^{r-1}\Omega^{p+1}(U;F)\)，使得

\[
(d \alpha) _ {x} (t _ {0}, \dots , t _ {p}) = \sum_ {i = 0} ^ {p} (- 1) ^ {i} (d _ {x} \tilde {\alpha}) (t _ {i}) (t _ {0}, \dots , \hat {t} _ {i}, \dots , t _ {p}) ^ {(1)}\tag{5}
\]

对所有 \(x \in U\) 以及 E 中的 \(t_{0}, \ldots, t_{p}\) 都成立。

8.3.5. 设 \(p\) 是一个 \(\geqslant 0\) 的整数，\(r\) 是 \(\mathbf{N}_{\mathbf{K}}\) 的一个元素。设 \(X\) 是 \(\mathbf{C}^{r+1}\) 类流形，\(F\) 是 Banach 空间，且 \(\omega \in {}^r\Omega^p(X;F)\)。存在且仅存在一个微分形式 \(\pi \in {}^{r-1}\Omega^{p+1}(X;F)\)，使得对于每个坐标图 \(c = (U,\varphi,E)\)（它是 \(X\) 的坐标图），若 \(\omega_c\) 是 \(\varphi(U)\) 上满足 \(\omega|U = \varphi^*(\omega_c)\) 的微分形式，则 \(\pi_c = d\omega_c\)，其中 \(d\omega_c\) 由 8.3.4 的公式 (5) 定义。

微分形式 \(\pi\) 称为 \(\omega\) 的外微分，记为 \(d\omega\)。它具有下列性质：

\begin{enumerate}
\def\labelenumi{(\arabic{enumi})}
\setcounter{enumi}{5}
\item
  若 \(\psi: Y \to X\) 是 \(C^{r+1}\) 类簇的态射且 \(\omega \in {}^r\Omega^n(X; F)\)，则有 \(\psi^*(d\omega) = d(\psi^*\omega)\)；特别地，\(d\) 与限制到子流形的运算可交换。
\item
  当 \(p = 0\) 时，映射 \(d: ^r\Omega^0(\mathbf{X}; \mathbf{F}) \to {}^{r-1}\Omega^1(\mathbf{X}; \mathbf{F})\) 与 8.2.2 中定义的映射一致。
\item
  若 \(r \geqslant 2\) 且 \(\omega \in {}^r\Omega^p(\mathbf{X};\mathbf{F})\)，则有 \(d(d\omega) = 0\)。\(^{(2)}\)
\item
  对 Banach 空间的每个连续线性映射 \(u: \mathbf{F} \to \mathbf{F}'\)，对于 \(d(u(\omega)) = u(d(\omega))\)，都有 \(\omega \in {}^{r}\Omega^{p}(\mathbf{U}; \mathbf{F})\)。
\item
  对 Banach 空间的每个连续双线性映射 \(F' \times F'' \rightarrow F\)，都有
\end{enumerate}

\[
d \left(\omega^ {\prime} \wedge \omega^ {\prime \prime}\right) = \left(d \omega^ {\prime}\right) \wedge \omega^ {\prime \prime} + (- 1) ^ {p ^ {\prime}} \omega^ {\prime} \wedge d \omega^ {\prime \prime}
\]

对于 \(\omega' \in {}^r\Omega^{p'}(\mathbf{X}; \mathbf{F}')\) 和 \(\omega'' \in {}^r\Omega^{p''}(\mathbf{X}; \mathbf{F}'')\)。


8.3.6. 设 \((u^{1},\ldots ,u^{n})\) 是 U 中 X 的一个坐标系。令

\[
\omega = \sum_ {i _ {1} <   \dots <   i _ {p}} f _ {i _ {1}, \dots , i _ {p}}. d u ^ {i _ {1}} \wedge \dots \wedge d u ^ {i _ {p}}
\]

\(^{1}\) 符号 \^{} 表示应省略它上方的符号（参见~7.8.4）。

\(^{2}\) 若 \(\omega\in^{1}\Omega^{p}(X;F)\) 且 \(d\omega\in^{1}\Omega^{p+1}(X;F)\)，仍有 \(d(d\omega)=0\)。

是 U 上的微分形式（其中 \(f_{i_1,\ldots,i_p} \in \mathcal{C}^r(\mathrm{U};\mathrm{F})\)）。于是有

\[
d \omega = \sum_ {i _ {1} <   \dots <   i _ {p}} d f _ {i _ {1}, \dots , i _ {p}} \wedge d u ^ {i _ {1}} \wedge \dots \wedge d u ^ {i _ {p}}.\tag{11}
\]

8.3.7. 假设 K 的特征为零。设 \(\omega \in {}^r\Omega^p (\mathbf{X};\mathbf{F})\) 是一个次数 \(p\geqslant 1\) 的微分形式，满足 \(d\omega = 0\)。对于每个 \(x\in \mathbf{X}\)，存在 \(x\) 的一个开邻域 U 以及一个 \(\pi \in {}^r\Omega^{p - 1}(\mathrm{U};\mathrm{F})\) 中的微分形式，使得在 U 上 \(d\pi = \omega\)。

